%===================================================================%
\section{Cloth Simulation in Robotics}
\label{sec:cloth_sim_robotics}
%===================================================================%

Robotic manipulation of textiles and garments requires reasoning over objects with
high-dimensional, non-linear dynamics, frequent self-contact, and pronounced sensitivity to
material parameters (e.g., bending stiffness, friction, and damping). These properties break
common assumptions used in rigid-body manipulation: applied forces not only translate the
object but also change its shape, topology-relevant features (e.g., folds, holes), and contact
state in a history-dependent manner. Such effects are central in tasks like folding, flattening,
dressing assistance, and sorting, and are a key contributor to the ``reality gap'' between
simulated and real cloth behavior \cite{longhini2025review_cloth_manipulation}.
Consequently, robotics simulators typically trade physical fidelity for robustness and speed,
favoring models and solvers that remain stable under large time steps and frequent contacts.

\subsection{Modeling approaches}
Common cloth models include (i) \emph{mesh-based mass--spring(-damper)} systems,
(ii) \emph{particle-based constraint} methods, and (iii) \emph{continuum} formulations
(e.g., \ac{FEM}/\ac{MPM}). Mass--spring models are intuitive but require careful tuning and may
struggle with complex elastic effects. Continuum models offer parameters with clearer
physical meaning, but are often computationally expensive for interactive robotics workloads.
Particle-based methods represent cloth as particles connected via holonomic constraints and
are widely used in real-time simulation due to their speed, stability, and ability to model
inextensibility and collisions \cite{longhini2025review_cloth_manipulation}.
In the context of this thesis, Isaac~Sim~5.1 provides cloth-capable simulation pathways in its
physics stack \cite{Muster.2017}.

\subsection{Position-Based Dynamics} %  for particle cloth

\ac{PBD} represents a deformable object by particle positions
$x_i \in \mathbb{R}^3$ with masses $m_i$ and enforces behavior through a set of constraints
$C_j(\mathbf{x})$ (equality/inequality), each with a stiffness parameter $k_j \in [0,1]$
\cite{muller2007pbd}. A typical \ac{PBD} time step (schematically) performs:
(i) external-force velocity update, (ii) position prediction (explicit step),
(iii) iterative constraint projection (Gauss--Seidel style), and (iv) velocity update from the
corrected positions \cite{muller2007pbd}. Let $\mathbf{p}$ be the predicted positions; the
solver iterates over constraints and applies local corrections to satisfy them. Linearization of
a constraint $C_j$ yields a correction of the form
\begin{equation}
\Delta \mathbf{x} = k_j \cdot s_j \cdot \mathbf{M}^{-1}\nabla C_j(\mathbf{x}), \qquad
s_j = \frac{-C_j(\mathbf{x})}{\nabla C_j(\mathbf{x})^\top \mathbf{M}^{-1} \nabla C_j(\mathbf{x})},
\label{eq:pbd_delta}
\end{equation}
which is applied repeatedly in a Gauss-Seidel fashion across all constraints
\cite{macklin2016xpbd,muller2007pbd}. This local projection viewpoint explains two
practically important properties: (a) \emph{robustness} under contacts and large constraint
sets, and (b) \emph{stiffness dependence} on time step and iteration count, i.e., constraints
become effectively stiffer with more solver iterations (and smaller $\Delta t$)
\cite{macklin2016xpbd,muller2007pbd}.
\ac{XPBD} extends \ac{PBD} by introducing a compliance formulation with Lagrange multipliers,
making stiffness less dependent on time step and iteration count while enabling force
estimates \cite{macklin2016xpbd}. While \ac{XPBD} is not strictly required for particle cloth, it
is relevant when aligning simulation parameters with physically meaningful stiffness values.

% \subsection{Why PhysX (Isaac~Sim) uses PBD for particle/cloth simulation}

The PhysX particle system uses \ac{PBD} to solve particle-to-particle dynamics and can represent
a range of phenomena, including deformables and cloth-like behavior \cite{physx_pbd_particlesystem_510,physx_particlesystem_513}.
In Omniverse/Isaac, ``particle cloth'' is explicitly described as a \emph{\ac{PBD} mass--spring}
system, distinct from newer surface/volume deformables that use different underlying models
\cite{omni_deformable_migration_guide}. The Omni Physics particles documentation further
emphasizes \ac{GPU}-accelerated \ac{PBD} particles supporting \emph{fluids, granular media, and
cloth}, including interaction with rigid bodies and articulations—an important requirement
for robotics scenes \cite{omni_physics_particles}.
From an algorithmic perspective, \ac{PBD} maps well to \ac{GPU} parallelism because constraints
can be partitioned and solved in parallel; e.g., a common approach constructs a graph over
constraints and applies greedy graph coloring so each color class can be processed
concurrently on the \ac{GPU} \cite{fratarcangeli2013gpu_pbd}. This aligns with Isaac~Sim's focus
on high-throughput simulation for learning workflows, and contributes to why \ac{PBD}-based
particle cloth has been a practical choice in PhysX-based pipelines.

\subsection{Vertex Block Descent (VBD)} % implicit, energy-based solver

\ac{VBD} targets the variational form of implicit Euler time integration.
Given positions $\mathbf{x}_t$ and velocities $\mathbf{v}_t$, the next-step positions are
obtained by minimizing a global energy
\begin{equation}
\mathbf{x}_{t+1} = \arg\min_{\mathbf{x}}\, G(\mathbf{x}), \qquad
G(\mathbf{x}) = \frac{1}{2h^2}\|\mathbf{x}-\mathbf{y}\|_{\mathbf{M}}^2 + E(\mathbf{x}),
\label{eq:vbd_global_energy}
\end{equation}
where $h$ is the time step, $\mathbf{y}=\mathbf{x}_t+h\mathbf{v}_t+h^2\mathbf{a}_\mathrm{ext}$,
and $E(\mathbf{x})$ is the total potential energy from internal forces and constraints
\cite{chen2024vbd}. \ac{VBD} solves this optimization by \emph{vertex-level block coordinate
descent}: iterating over vertices and updating one vertex position at a time by minimizing a
local energy $G_i(\mathbf{x})$ that depends only on that vertex and its incident force elements
\cite{chen2024vbd}. Concretely, for vertex $i$:
\begin{equation}
\mathbf{x}_i \leftarrow \arg\min_{\mathbf{x}_i} G_i(\mathbf{x}), \qquad
G_i(\mathbf{x})=\frac{m_i}{2h^2}\|\mathbf{x}_i-\mathbf{y}_i\|^2+\sum_{j\in F_i}E_j(\mathbf{x}),
\label{eq:vbd_local_energy}
\end{equation}
and a practical update uses a small $3\times 3$ Newton step
\begin{equation}
\mathbf{H}_i \Delta \mathbf{x}_i = \mathbf{f}_i,
\label{eq:vbd_newton}
\end{equation}
with $\mathbf{H}_i$ the local Hessian and $\mathbf{f}_i$ the local force/gradient term
\cite{chen2024vbd}. After updating positions, velocities follow from the implicit Euler
relationship $\mathbf{v}_{t+1}=(\mathbf{x}_{t+1}-\mathbf{x}_t)/h$ \cite{chen2024vbd}.
A key practical advantage is that \ac{VBD} is designed to achieve convergence to the implicit
Euler solution with unconditional stability while remaining highly parallelizable via vertex
coloring schemes \cite{chen2024vbd}.

% \subsection{Why newer simulation engines (Newton) employ VBD for cloth}

Newton exposes a \ac{VBD}-based solver for particles/thin deformables (and related variants for
rigid bodies), reflecting a design choice toward implicit, energy-based integration for robust
simulation of stiff systems \cite{newton_solvers_doc}. In NVIDIA's robotics context, Newton
is presented as enabling improved performance and stability of a \ac{VBD} solver for thin
deformables such as clothing, supporting large-scale, \ac{GPU}-accelerated learning workloads
\cite{isaaclab_newton_blog_2025}.
For cloth manipulation in robotics, these properties are attractive: cloth often mixes highly
stiff constraints (near-inextensible stretch) with softer bending, while undergoing frequent
contacts; implicit/variational formulations can better tolerate such stiffness ratios and larger
time steps without requiring excessive constraint iterations.

\subsection{Comparison} % PBD (PhysX particle cloth) vs.\ VBD (Newton thin deformables)

Both methods iterate locally and can be implemented in \ac{GPU}-friendly Gauss--Seidel-like
passes, but they differ in their mathematical target:
\begin{itemize}
  \item \textbf{\acs{PBD}} projects positions to satisfy constraints. It is simple, robust, and fast,
  but its effective stiffness depends on time step and iteration count, and mapping
  parameters to physical material properties can be difficult \cite{muller2007pbd,macklin2016xpbd,longhini2024review_cloth_manipulation}.
  \item \textbf{\acs{VBD}} directly minimizes the variational implicit Euler energy, providing a
  principled implicit integration target with strong stability guarantees and favorable
  convergence properties in large \ac{DoF} systems \cite{chen2024vbd}.
\end{itemize}
In Omniverse's physics ecosystem, this distinction is reflected at the feature level: particle
cloth is characterized as a \ac{PBD} mass--spring system, while newer deformable features
(and adjacent solvers) are moving toward implicit, energy/compliance-based formulations
\cite{omni_deformable_migration_guide}. For this thesis, this motivates using PhysX/\ac{PBD}
particle cloth when prioritizing throughput and robustness in Isaac~Sim~5.1, while also
motivating Newton/\ac{VBD} when the task demands greater stability under stiff cloth behavior
and larger time steps in learning-centric simulation pipelines.

% ---------------------------------------------------------------------------
% BibTeX keys:
%  - longhini2024review_cloth_manipulation  (cloth_manipulation_review.pdf)
%  - muller2007pbd                           (pbd_paper.pdf)
%  - macklin2016xpbd                         (xpbd_paper.pdf)
%  - fratarcangeli2013gpu_pbd                (gpu_pbd_paper.pdf)
%  - chen2024vbd                             (vbd_paper.pdf)
%  - physx_pbd_particlesystem_510 / _513      (PhysX docs pages)
%  - omni_physics_particles                  (Omni Physics particles docs page)
%  - omni_deformable_migration_guide         (Omni Physics deformables migration guide)
%  - newton_solvers_doc                      (Newton solver docs page)
%  - isaaclab_newton_blog_2025               (NVIDIA developer blog post)
% ---------------------------------------------------------------------------

% Web citations:
% \footnote{\url{https://nvidia-omniverse.github.io/PhysX/physx/5.1.0/_build/physx/latest/class_px_physics.html}}
% \footnote{\url{https://nvidia-omniverse.github.io/PhysX/physx/5.1.3/docs/ParticleSystem.html}}
% \footnote{\url{https://docs.omniverse.nvidia.com/kit/docs/omni_physics/latest/dev_guide/particles/particles.html}}
% \footnote{\url{https://docs.omniverse.nvidia.com/kit/docs/omni_physics/107.3/dev_guide/deformables_beta/deformable_migration.html}}
% \footnote{\url{https://newton-physics.github.io/newton/api/newton_solvers.html}}
% \footnote{\url{https://developer.nvidia.com/blog/train-a-quadruped-locomotion-policy-and-simulate-cloth-manipulation-with-nvidia-isaac-lab-and-newton/}}
